\section{Econometric Application}
\subsection{Targeted use case}
The proposed NS-SDN architecture is designed for macroeconomic and financial time series that exhibit secular trend and evolving cycle durations and volatility \citep{hamilton1989,lubik2015}, where the amplitude, frequency, and phase of cycles may evolve over time due to regime changes, policy shocks, or structural breaks.

Canonical examples include: real GDP and output gap measures; inflation and inflation expectations; interest rates and yield spreads; unemployment and labor market indicators; and commodity prices and real exchange rates. These series are known to exhibit (i) persistent low-frequency trends, (ii) medium-frequency business-cycle dynamics, and (iii) episodic nonstationarity following shocks (e.g., monetary tightening, supply disruptions). Traditional linear constant-parameter models may fail to capture such evolving macroeconomic dynamics \citep{lubik2015}.

\subsection{Mapping NS-SDN to econometric structure}
NS-SDN can be interpreted as a nonlinear, nonstationary unobserved-components model, in which the latent state $h_n$ plays the role of an unobserved economic state; the trend component $B_n$ corresponds to a stochastic trend; the sinusoidal components correspond to stochastic cycles with time-varying amplitude and frequency; and the observation equation is explicitly additive and interpretable.

Unlike classical unobserved-components models or Kalman-filter-based stochastic cycle models, NS-SDN does not impose fixed frequencies, does not assume linear Gaussian dynamics, and allows regime-dependent evolution of spectral structure. This makes NS-SDN particularly suitable for environments where cyclical behavior itself is nonstationary, such as post-crisis macroeconomic regimes.

\subsection{Forecasting task formulation}
Let $\{y_1, \dots, y_T\}$ denote an observed time series. The forecasting task is defined as predicting future values $\hat y_{T+1}, \hat y_{T+2}, \dots, \hat y_{T+H}$ given information available up to time $T$.

NS-SDN is trained in a recursive one-step-ahead forecasting framework, where the latent state $h_n$ is updated sequentially, predictions are generated via the spectral emission equation, and multi-step forecasts are generated recursively, consistent with standard practice in neural autoregressive forecasting \citep{salinas2020}. This setup mirrors the forecasting protocols used for VARs, state-space models, and neural sequence models, enabling fair comparison.

